% ============================================================
% Capitulo 2 -- Fundamentacao Teorica
% Tamanho-alvo: 12-18 paginas | Sem figuras (visualizacao vai pro Cap 4)
% Ordem de escrita: PENULTIMO (depois de Resultados)
% ============================================================

\chapter{Fundamenta\c{c}\~ao Te\'orica}
\label{cap:fundamentacao}

% ============================================================
\section{Detec\c{c}\~ao de fake news: panorama}
\label{sec:panorama}

% TESE: abordagens isoladas (NLP-only ou estrutura-only) tem limites;
%       trabalhos recentes movem para hibridizacao.
%
% CONTEUDO (~1.5 pagina):
% - NLP classico em deteccao (BoW, TF-IDF, LSTM, BERT).
% - Limites observados em LLMs gerando texto convincente.
% - Direcao hibrida (texto + contexto + propagacao).
%
% CITACOES: \cite{pmc_overviewfakenews}, \cite{ieee_fakenewslandscape},
%           \cite{mdpi_cnn_llm_nlp}.

% ============================================================
\section{Grafos e redes sociais}
\label{sec:grafos}

% TESE: propagacao em rede social = cascata enraizada (arvore com raiz).
%
% CONTEUDO (~1 pagina):
% - Definicoes basicas: grafo dirigido, arvore, raiz, profundidade,
%   branching factor.
% - Propagacao como cascata: post original = raiz; reposts = filhos.
%
% CITACAO opcional: \cite{newman2003structure}.

% ============================================================
\section{Embeddings textuais}
\label{sec:embeddings}

% TESE: usamos Sentence-BERT multilingue (paraphrase-multilingual-mpnet-base-v2, 768 dim),
%       NAO BERT vanilla.
%
% CONTEUDO (~1 pagina):
% - word embeddings -> BERT -> Sentence-BERT.
% - Modelo usado: paraphrase-multilingual-mpnet-base-v2 (768 dim).
% - Pegadinha: Sentence-BERT != BERT vanilla.
%
% CITACOES: \cite{devlin2019bert}, \cite{reimers2019sbert}.

% ============================================================
\section{Redes Neurais de Grafos}
\label{sec:gnns}

% TESE: GNN = aprendizado de representacao via agregacao de vizinhos.
%       Apresentar as 3 familias usadas (GCN, GAT, SAGE).

\subsection{Graph Convolutional Networks (GCN)}
\label{sec:gcn}

% TESE: agregacao espectral com normalizacao simetrica por grau.
%
% CONTEUDO: regra de propagacao + intuicao + complexidade.

\begin{equation}
H^{(l+1)} = \sigma\left(\hat{D}^{-1/2}\hat{A}\hat{D}^{-1/2} H^{(l)} W^{(l)}\right)
\label{eq:gcn}
\end{equation}

% CITACAO: \cite{kipf2017gcn}.

\subsection{Graph Attention Networks (GAT)}
\label{sec:gat}

% TESE: introduz coeficiente de atencao aprendido por aresta.
%
% CONTEUDO: equacao de attention + multi-head.

\begin{equation}
\alpha_{ij} = \frac{\exp\left(\mathrm{LeakyReLU}(\mathbf{a}^T [W h_i \,\|\, W h_j])\right)}{\sum_{k \in \mathcal{N}(i)} \exp\left(\mathrm{LeakyReLU}(\mathbf{a}^T [W h_i \,\|\, W h_k])\right)}
\label{eq:gat_alpha}
\end{equation}

% CITACAO: \cite{velickovic2018gat}.

\subsection{GraphSAGE}
\label{sec:sage}

% TESE: agregacao indutiva (amostra vizinhanca + agrega), generaliza
%       para nos novos.
%
% CONTEUDO:
% - Agregadores: mean, max, lstm (usamos mean).
% - Pooling final: global_mean_pool.
% - Implementacao via PyTorch Geometric.
%
% CITACOES: \cite{hamilton2017graphsage}, \cite{fey2019pyg}.

% ============================================================
\section{Degenera\c{c}\~ao do GCN com features homog\^eneas}
\label{sec:degeneracao}

% TESE CRITICA: se X_i = X_j para todo i,j, GCN reduz a transformacao
%               escalar dependente apenas do grau -- sem poder discriminativo.
%               Justifica adicionar encodings posicionais.
%
% CONTEUDO (~1.5 pagina):
% - Notacao: X em R^{NxD}, \tilde{A} = D^{-1/2}\hat{A}D^{-1/2}.
% - Cenario: X = 1_N x^T (toda linha igual).
% - Demonstracao matematica (equacao abaixo).
% - Conclusao: h_i^{(1)} e' multiplo escalar do mesmo xW, fator = grau de i.

\begin{align}
H^{(1)} &= \sigma\left(\tilde{A} X W\right) \notag \\
        &= \sigma\left(\tilde{A} \mathbf{1}_N x^T W\right) \notag \\
        &= \sigma\left((\tilde{A} \mathbf{1}_N)(x W)^T\right)
\label{eq:degeneracao}
\end{align}

% IMPLICACAO: justifica \S\ref{sec:encodings}.

% ============================================================
\section{Encodings posicionais}
\label{sec:encodings}

% TESE: 3 features adicionadas por no quebram a degenerencia de \S\ref{sec:degeneracao}.
%
% CONTEUDO (~1 pagina):
% - is_root \in {0,1}: 1 se no e' a raiz da arvore.
% - grau_norm = deg(v) / deg_max_global \in [0,1].
% - pos = ordem BFS / num_filhos, normalizado em [0,1].
% - Como cada um quebra a degenerencia.
%
% CITACAO opcional: \cite{arxiv_topologicalfeatures}.

% ============================================================
\section{Explicabilidade em GNNs: GNNExplainer}
\label{sec:gnn_explainer}

% TESE: GNNExplainer aprende mascara de arestas/features maximizando
%       mutual info com a predicao.
%
% CONTEUDO (~1 pagina):
% - Problema de otimizacao (MI entre predicao e subgrafo mascarado).
% - Setup pratico: torch_geometric.explain.Explainer com epochs=200.
%
% CITACAO: \cite{ying2019gnnexplainer}.

% ============================================================
\section{Valida\c{c}\~ao estat\'istica}
\label{sec:estatistica}

% TESE: k-fold pareado + ttest_rel + Cohen's d sao a base estatistica.
%       Bootstrap pareado para subsets pos-hoc (Cap 4).
%
% CONTEUDO (~1 pagina):
% - K-fold estratificado: por que estratificado, por que 10 folds.
% - ttest_rel (pareado) vs ttest_ind: por que pareado tem mais poder.
% - Cohen's d como effect size; tabela de interpretacao abaixo.
% - Bootstrap pareado: usado em subsets nao-pareados entre folds (Cap 4.8).
%
% CITACOES: \cite{cohen1988statistical}, \cite{efron1993bootstrap}.

\begin{table}[h]\centering
\caption{Interpreta\c{c}\~ao de tamanho de efeito de Cohen ($d$).}
\label{tab:cohen_d_interpretacao}
\begin{tabular}{ll}
\toprule
$|d|$           & Interpreta\c{c}\~ao \\
\midrule
$< 0.2$         & desprez\'ivel \\
$0.2 \leq |d| < 0.5$ & pequeno \\
$0.5 \leq |d| < 0.8$ & m\'edio \\
$\geq 0.8$      & grande \\
\bottomrule
\end{tabular}
\end{table}

% ============================================================
\section{Trabalhos relacionados}
\label{sec:relacionados}

% TESE: posicionar este trabalho como diagnostico do paradigma topologico,
%       nao como nova arquitetura.

\subsection{Surveys de GNNs em fake news}
\label{sec:surveys}
% CITACOES: \cite{pmc_gnnsurveyfakenews}, \cite{pmc_overviewfakenews}.

\subsection{M\'etodos GNN can\^onicos}
\label{sec:metodos_canonicos}
% CITACOES: \cite{monti2019fakenews}, \cite{dou2021upfd}, \cite{bian2020bigcn},
%           \cite{mdpi_bidirectionalgcn}.

\subsection{Detec\c{c}\~ao multil\'ingue}
\label{sec:multilingue}
% CITACOES: \cite{acl_hemt}, \cite{crosslanguagefakenews}, \cite{pmc_multilangs_covid}.

\subsection{Trabalhos brasileiros relacionados}
\label{sec:trabalhos_br}
% CITACOES: \cite{ufc_fakewhatsapp}, \cite{pucsp_gnnfakenews}, \cite{usp_fakenews}.

\subsection{Homofilia e assortatividade em redes}
\label{sec:homofilia}

% TESE: GNNs convolucionais dependem de assortatividade por classe.
%       Fundamenta teoricamente por que GossipCop funciona (Cap 4.5).
%
% CITACOES: \cite{newman2003mixing}, \cite{pei2020geomgcn}.

\subsection{Posicionamento deste trabalho}
\label{sec:posicionamento}

% TESE: nao propomos nova arquitetura; propomos diagnostico de
%       quando o paradigma funciona.
%
% CONTEUDO (~1 paragrafo):
% - Reforcar que comparacao com BiGCN/GCNFN e' via valores publicados
%   (limitacao em \S\ref{sec:limitacoes}).
